\section{Énoncés des Exercices}
\subsection{Séries numériques}

% --- (Source: EXERCICE 7 analyse) ---
\begin{exercice}[Exercice 7 Analyse]
\begin{enumerate}
    \item Soient $(u_{n})$ et $(v_{n})$ deux suites de nombres réels positifs.
        On suppose que $(v_{n})$ est non nulle à partir d'un certain rang.
        
        Montrer que si $u_n \sim v_n$ alors $\sum u_n$ et $\sum v_n$ sont de même nature.

    \item Étudier la convergence de la série:
        \[\sum_{n\geqslant2}\frac{(i-1)\ln n \sin\left(\frac{1}{n}\right)}{(\sqrt{n+3}-1)}\]
        (où $i$ est le complexe tel que $i^2 = -1$).
\end{enumerate}
\end{exercice}

% --- (Source: EXERCICE 46 analyse) ---
\begin{exercice}[Exercice 46 Analyse]
    On considère la série $\sum_{n\geqslant1}\cos\left(\pi\sqrt{n^{2}+n+1}\right)$.
    \begin{enumerate}
        \item Prouver que, au voisinage de $+\infty$, $\pi\sqrt{n^{2}+n+1}=n\pi+\frac{\pi}{2}+\alpha\frac{\pi}{n}+O\left(\frac{1}{n^{2}}\right)$ 
            où $\alpha$ est un réel que l'on déterminera.
        \item En déduire que $\sum_{n\geqslant1}\cos\left(\pi\sqrt{n^{2}+n+1}\right)$ converge.
        \item La série converge-t-elle absolument ?
    \end{enumerate}
\end{exercice}

% --- (Source: Exercice 1) ---
\begin{exercice}
\begin{enumerate}
    \item Pour $s\in\C$ tel que $\text{Re}(s)>1$, justifier la convergence absolue de la série de Riemann de somme :
        \[\zeta(s)=\sum_{n=1}^{+ \infty}\frac{1}{n^{s}}\]

    \item Pour $s\in\R$ tel que $s>0$, justifier la convergence de la série alternée de Riemann de somme :
        \[A(s)=\sum_{n=1}^{+ \infty}\frac{(-1)^{n}}{n^{s}}\]
    
    \item  Pour $s\in\C$ tel que $\text{Re}(s)>1$, simplifier $A(s)+\zeta(s)$ et en déduire un lien entre $\zeta(s)$ et $A(s)$.
    \item On s'intéresse à la convergence de $A(s)$ pour $s$ complexe tel que $\text{Re}(s)>0$.
        On fixe $s=a+ib$ avec $a>0$ et $b\in\R$ et on pose $u_{n}=\frac{(-1)^{n}}{n^{s}}$. 
        On note $v_{n}=\text{Re}(u_{n})$ et $w_{n}=\text{Im}(u_{n})$.
    \begin{enumerate}
        \item On pose $V_{n}=v_{2n}+v_{2n+1}$. En montrant que $\frac{1}{(2n+1)^{a}}=\frac{1}{(2n)^{a}}+O\left(\frac{1}{n^{1+a}}\right)$, 
            montrer que la série $\sum V_{n}$ est absolument convergente.
        \item En déduire que la série $\sum v_{n}$ converge.
        \item Conclure en raisonnant de même avec $\sum w_{n}$.
    \end{enumerate}
\end{enumerate}
\textit{Note : Il existait une question 5 qui évoquait de prouver l'hypothèse de Riemann, elle n'est pas traitée ici.}
\end{exercice}

% --- (Source: Exercice 2) ---
\begin{exercice}
\begin{enumerate}
    \item
    \begin{enumerate}
        \item Calculer pour $x\in[0,1]$, $\displaystyle\sum_{k=0}^{N}\left(-x^{2}\right)^{k}$.
        \item Intégrer cette égalité sur $[0, 1]$.
        \item En déduire que $\displaystyle\sum_{n=0}^{+\infty}\frac{(-1)^{n}}{2n+1}=\frac{\pi}{4}$.

    \end{enumerate}
    
    \item On définit pour tout $n$, $\displaystyle I_{n}=\int_{0}^{\pi/2}x^{2}\frac{\sin(2nx)}{\sin x} \dx$.
    \begin{enumerate}
        \item Montrer que $\displaystyle I_{n+1}-I_{n}=\int_{0}^{\pi/2}2x^{2}\cos((2n+1)x) \dx$ 
        et en déduire que $I_{n+1}-I_{n}=\frac{(-1)^{n}\pi^{2}}{2(2n+1)}-\frac{4(-1)^{n}}{(2n+1)^{3}}$.
        \item On pose pour $x>0$, $f(x)=\frac{x^{2}}{\sin x}$. Montrer que $f$ se prolonge en 0 en une fonction de classe $\mathscr C^{1}$.
        \item Utiliser une intégration par parties pour montrer que $\displaystyle|I_{n}|\leqslant\frac{|f(\pi/2)|+\int_{0}^{\pi/2}|f^{\prime}(x)| \dx}{2n}$.
        \item En déduire que $\displaystyle\sum_{n=0}^{+\infty}\frac{(-1)^{n}}{(2n+1)^{3}}=\frac{\pi^{3}}{32}$. % Note: PDF says pi^3/32
    \end{enumerate}
\end{enumerate}
\end{exercice}

% --- (Source: Exercice 3) ---
\begin{exercice}
    On pose pour $\theta\in]0,\pi[$, $u_{n}=\frac{\cos n\theta}{n}$.
    \begin{enumerate}
        \item On pose $\displaystyle C_{n}=\sum_{k=0}^{n}\cos(k\theta)$. Montrer que cette suite est bornée.
        \item En remarquant que pour tout $n\geqslant1$, $\cos n\theta=C_{n}-C_{n-1}$ et en considérant sa somme partielle, 
            montrer que la série $\sum u_{n}$ converge.
        \item Montrer que $|\cos(n\theta)|\geqslant \cos^{2}(n\theta)$ et en déduire que la série $\sum u_{n}$ ne converge pas absolument.
    \end{enumerate}
\end{exercice}

% --- (Source: Exercice 4) ---
\begin{exercice}
    Préciser la nature des séries suivantes:
    \[\sum\frac{(-1)^{n}}{3n-1}, \quad \sum\frac{(-1)^{n}}{n^{\alpha}-(-1)^{n}n^{\beta}} \ (\alpha\ne\beta), \quad \sum \ln\left(1+(-1)^{n}\frac{\ln n}{n^{\alpha}}\right) \ (\alpha>0)\]
\end{exercice}

% --- (Source: Exercice 5 - Stirling) ---
\begin{exercice}
    Montrer que pour tout naturel $n$, $n!e=p_{n}+\frac{1}{n+1}+O\left(\frac{1}{n^{2}}\right)$ 
    où $p_{n}$ est un entier dont on précisera la parité en fonction de $n$.
    En déduire la nature des séries $\sum \sin(n!2\pi e)$ et $\sum \sin(n!\pi e)$.
\end{exercice}

% --- (Source: Exercice 6 - Cesaro) ---
\begin{exercice}
Soit $(\sum u_{n})$ une série numérique convergente. On pose $\displaystyle S_{n}=\sum_{k=0}^{n}u_{k}$.
\begin{enumerate}
    \item Exprimer $\displaystyle\sum_{k=1}^{n}ku_{k}$ à l'aide de $\displaystyle\sum_{k=1}^{n}S_{k}$ (utiliser $u_n = S_n - S_{n-1}$).
    \item En déduire que $\displaystyle\lim_{n\rightarrow+\infty}\frac{1}{n}\sum_{k=1}^{n}ku_{k}=0$.
    \item On pose $\displaystyle v_{n}=\frac{1}{n(n+1)}\sum_{k=1}^{n}ku_{k}$.
        Exprimer $\displaystyle\sum_{k=1}^{n}v_{k}$ à l'aide de $S_{n}$.
        En déduire que la série $(\sum v_{n})$ converge et a même somme que $(\sum u_{n})$.
\end{enumerate}
\end{exercice}

\begin{exercice}
Montrer l'existence et calculer : \[\sum_{n=0}^{+\infty} \frac{1}{(3n)!}\]
\end{exercice}
    
\subsection{Séries dans un espace vectoriel normé de dimension finie}

% --- (Source: EXERCICE 40 analyse) ---
\begin{exercice}[Exercice 40 Analyse]
    Soit $\mathcal{A}$ une algèbre de dimension finie admettant $e$ pour élément unité et munie d'une norme telle que $\|uv\| \leqslant \|u\|\|v\|$.
    \begin{enumerate}
        \item Soit $u \in \mathcal{A}$ tel que $\|u\| < 1$.
        \begin{enumerate}
            \item Démontrer que la série $\sum_{n\geqslant1} u^{n}$ est convergente.
            \item Démontrer que $(e-u)$ est inversible et que $\displaystyle (e-u)^{-1}=\sum_{n=0}^{+\infty}u^{n}$.
        \end{enumerate}
        \item Démontrer que, pour tout $u \in \mathcal{A}$, la série $\sum\frac{u^{n}}{n!}$ converge.
    \end{enumerate}
\end{exercice}

% --- (Source: EXERCICE 61 algèbre) ---
\begin{exercice}[Exercice 61 Algèbre]
    On note $\mathcal{M}_{n}(\C)$ l'espace des matrices carrées d'ordre $n$. Pour $A=(a_{i,j})$, on pose $\|A\| = \max_{i,j \in \ninter 1 n} |a_{i,j}|$.
    \begin{enumerate}
        \item Prouver que $\|\cdot\|$ est une norme.
        \item Démontrer que $\|AB\| \leqslant n\|A\|\|B\|$ et $\forall p \in \N^* \|A^p\| \leqslant n^{p-1}\|A\|^p$.
        \item Démontrer que la série $\sum\frac{A^{p}}{p!}$ est absolument convergente. Est-elle convergente ?
    \end{enumerate}
\end{exercice}

% --- (Source: Exercice 7) ---
\begin{exercice}
    Calculer les exponentielles des matrices suivantes:
    \[\begin{pmatrix} 0 & a \\ -a & 0 \end{pmatrix}, \quad \begin{pmatrix} 1 & 1 \\ 0 & 1 \end{pmatrix}\]

\end{exercice}

% --- (Source: Exercice 8) ---
\begin{exercice}
    On pose $A=\begin{pmatrix}3&6\\ 4&1\end{pmatrix}$.
    \begin{enumerate}
        \item Déterminer deux suites $(a_{n})$ et $(b_{n})$ telles que $A^{n}=a_{n}(A-7I_{2})+b_{n}(A+3I_{2})$.
        \item En déduire $\exp(A)$.
    \end{enumerate}
\end{exercice}

% --- (Source: Exercice 9) ---
\begin{exercice}
    Montrer que pour une matrice $A\in \mathcal{M}_{p}(\K)$ de norme suffisamment petite la série $\sum_{n\geqslant0}A^{n}$ est absolument convergente. 
    Montrer que sa somme est l'inverse de $I_{p}-A$.
\end{exercice}

% --- (Source: Exercice 10) ---
\begin{exercice}
    On munit $\mathcal{M}_{p}(\C)$ d'une norme d'algèbre.
    \begin{enumerate}
        \item Pour tout naturel $n$ et tout $k\in\{0,\dots,n\}$ on pose $a_{n,k}=\frac{1}{k!}-\frac{n!}{k!(n-k)!n^{k}}$. 
            Montrer que $a_{n,k}\geqslant0$.
        \item On pose $M_{n}=(I_{p}+\frac{1}{n}A)^{n}$.
        \begin{enumerate}
            \item Montrer que $\displaystyle M_{n}-\sum_{k=0}^{n}\frac{A^k}{k!}$ tend vers 0.
            \item En déduire que la suite $(M_{n})$ converge vers $\exp(A)$.
        \end{enumerate}
    \end{enumerate}
\end{exercice}

\subsection{Séries à termes positifs}

% --- (Source: EXERCICE 5 analyse) ---
\begin{exercice}[Exercice 5 Analyse]
\begin{enumerate}
    \item On considère la série de terme général $u_{n}=\frac{1}{n(\ln n)^{\alpha}}$ ($n\geqslant2, \alpha\in\R$).
    \begin{itemize}
        \item Cas $\alpha\leqslant0$ : En utilisant une minoration simple, démontrer que la série diverge.
        \item Cas $\alpha>0$ : Étudier la nature de la série (comparaison série-intégrale avec $f(x)=\frac{1}{x(\ln x)^{\alpha}}$).
    \end{itemize}
    \item Déterminer la nature de la série $\sum_{n\geqslant3}\frac{\left(e-\left(1+\frac{1}{n}\right)^{n}\right)e^{\frac{1}{n}}}{(\ln(n^{2}+n))^{2}}$.
\end{enumerate}
\end{exercice}

% --- (Source: Exercice 11) ---
\begin{exercice}
    Établir la convergence et calculer la somme de la série télescopique : $\sum_{n\geqslant2}\ln\left(1-\frac{1}{n^{2}}\right)$.
\end{exercice}

% --- (Source: Exercice 12) ---
\begin{exercice}
\begin{enumerate}
    \item Comparer pour $a, b \in \R^{+}$, $\arctan(a) - \arctan(b)$ avec $\arctan \frac{a-b}{1+ab}$.
    \item En déduire la convergence et la somme de $\sum_{n\geqslant1}\arctan\left(\frac{1}{n^{2}+n+1}\right)$.
    \item On définit la suite de Fibonacci $(F_{n})$ par $F_{0}=0, F_{1}=1, F_{n+2}=F_{n}+F_{n+1}$.
    \begin{enumerate}
        \item Montrer la convergence de $\sum_{k\geqslant1}\arctan\left(\frac{1}{F_{2k+1}}\right)$.
        \item Montrer que pour tout $n$, $\displaystyle\frac{\pi}{4}=\sum_{k=1}^{n}\arctan\left(\frac{1}{F_{2k+1}}\right)+\arctan\left(\frac{1}{F_{2n+2}}\right)$.
        \item En déduire que $\displaystyle\sum_{k=1}^{+\infty}\arctan\left(\frac{1}{F_{2k+1}}\right) = \frac{\pi}{4}$.
    \end{enumerate}
\end{enumerate}
\end{exercice}

% --- (Source: Exercice 13) ---
\begin{exercice}
    Soient $a, b, c, d > 0$ et $u_{n}$ telle que $\frac{u_{n+1}}{u_{n}}=\frac{(n+a)(n+b)}{(n+c)(n+d)}$.
    \begin{enumerate}
        \item Pour $\alpha>0$, on pose $v_{n}=\ln(n^{\alpha}u_{n})$. Écrire un développement asymptotique de $v_{n+1}-v_{n}$. 
            À quelle condition la série télescopique converge-t-elle ?
        \item En déduire un équivalent de $u_{n}$.
    \end{enumerate}
\end{exercice}

% --- (Source: Exercice 14) ---
\begin{exercice}
    Si $\sum u_{n}$ converge (à termes positifs), établir la convergence des séries
    \[\sum u_{n}^{2},~ \sum\frac{u_{n}}{1+u_{n}^{2}}, ~\sum\frac{\sqrt{u_{n}}}{n^{\alpha}} ~\text{avec}~ \alpha>\frac{1}{2}\]
\end{exercice}

% --- (Source: analyse) ---
\begin{exercice}
    Soit $(u_{n})$ à termes strictement positifs.
    \begin{enumerate}
        \item Démontrer que si $\lim \frac{u_{n+1}}{u_{n}}=l < 1$, alors $\sum u_{n}$ converge
        \item Quelle est la nature de $\sum_{n\geqslant1}\frac{n!}{n^{n}}$ ?
    \end{enumerate}
\end{exercice}

% --- (Source: Exercice 15 - Raabe Duhamel) ---
\begin{exercice}[Règle de Raabe Duhamel]
\begin{enumerate}
    \item Montrer que si $\sum u_{n}$ et $\sum v_{n}$ sont deux séries à termes positifs 
        telles que pour $n$ assez grand $\frac{u_{n+1}}{u_{n}}\leqslant\frac{v_{n+1}}{v_{n}}$, 
        alors la convergence de $\sum v_{n}$ implique celle de $\sum u_{n}$
    \item On suppose $\frac{u_{n+1}}{u_{n}}=1-\frac{\alpha}{n}+o\left(\frac{1}{n}\right)$. 
        Montrer que $\sum u_{n}$ converge pour $\alpha>1$ et diverge pour $\alpha<1$ (comparer avec $v_n = \frac{1}{n^\beta}$).
\end{enumerate}
\end{exercice}

% --- (Source: Exercice 16) ---
\begin{exercice}
    Trouver la nature des séries :
    \[\sum \frac{1}{\ln(n)^{\ln n}}, \quad \sum \frac {n^n} {2^{n^2}}, \quad \sum \frac{n!^2}{(2n)!}, \quad \sum \left(\sin\frac{1}{n}\right) n^\alpha, \alpha \in \R\]
\end{exercice}

% --- (Source: Exercice 17) ---
\begin{exercice}
    Soit $\sum u_{n}$ une série à termes positifs, $u_0 > 0$. On pose $\displaystyle S_{n}=\sum_{k=0}^{n}u_{k}$.
    \begin{enumerate}
        \item Supposons $\sum u_{n}$ converge. Montrer que $\sum\frac{u_{n}}{S_{n}}$ converge
        \item Supposons $\sum\frac{u_{n}}{S_{n}}$ converge. 
            Trouver la nature de $\sum \ln\left(\frac{S_{n-1}}{S_{n}}\right)$. 
            En déduire que $\sum u_{n}$ converge
        \item Montrer que si $\sum a_{n}$ diverge (termes positifs), 
            il existe $\sum b_{n}$ divergente (à termes positifs) telle que $b_{n}=o(a_{n})$.
    \end{enumerate}
\end{exercice}

% --- (Source: Exercice 18) ---
\begin{exercice}
    Soit $\alpha>1$ et $\displaystyle R_{n}=\sum_{k=n+1}^{+\infty}\frac{1}{k^{\alpha}}$.
    \begin{enumerate}
        \item Déterminer un équivalent de $R_{n}$
        \item Condition nécessaire sur $\alpha$ de convergence de $\sum R_{n}$.
        \item Valeur de la somme dans ce cas ?
    \end{enumerate}
\end{exercice}

\newpage
\section{Corrections}

% --- Correction Exercice 2 (du fichier source, Exercice 46 ici) ---
\begin{correction}{2}
\begin{enumerate}
    \item On fait le calcul :
    \begin{align*}
        \pi\sqrt{n^2+n+1} &= \pi n \left(1 + \frac{1}{n} + \frac{1}{n^2}\right)^{1/2} \\
        &= \pi n \left(1 + \frac{1}{2n} + \frac{3}{8n^2} + O\left(\frac{1}{n^3}\right)\right) \\
        &= \pi n + \frac{\pi}{2} + \frac{3\pi}{8n} + O\left(\frac{1}{n^2}\right).
    \end{align*}
    
    \item $\cos(\pi\sqrt{\dots}) = \cos(n\pi + \pi/2 + \varepsilon_n) = (-1)^n \cos(\pi/2 + \varepsilon_n) = (-1)^{n+1} \sin(\varepsilon_n)$.
    Et $\sin(\varepsilon_n) \sim \frac{3\pi}{8n}$.

    C'est une série alternée $\sum (-1)^{n+1} \frac{K}{n}$. Elle converge (TSSA).
    Absolue convergence ? Non, car en $1/n$.
\end{enumerate}
\end{correction}

% --- Correction Exercice 3 ---
\begin{correction}{3}
\begin{enumerate}
    \item
        Soit $s=a+ib$ avec $a>1$.
        $|\frac{1}{n^s}| = \frac{1}{|n^a e^{ib\ln n}|} = \frac{1}{n^a}$.
        C'est une série de Riemann avec $a>1$, donc convergente. Absolue convergence vérifiée.

    \item 
        Pour $s>0$ réel, la suite $\frac{1}{n^s}$ tend vers 0 en décroissant.
        D'après le Théorème Spécial des Séries Alternées (TSSA), $\sum \frac{(-1)^n}{n^s}$ converge.

    \item 
        $A(s)+\zeta(s) = \sum_{n=1}^{+ \infty} \frac{(-1)^n+1}{n^s}$.
        Si $n$ est impair, $(-1)^n+1=0$. Si $n=2p$ est pair, $(-1)^n+1=2$.
        $$ A(s)+\zeta(s) = \sum_{p=1}^{+ \infty} \frac{2}{(2p)^s} = \frac{2}{2^s} \sum_{p=1}^{+ \infty} \frac{1}{p^s} = 2^{1-s} \zeta(s) $$
        D'où $A(s) = (2^{1-s}-1)\zeta(s)$.

    \item
    \begin{enumerate}
        \item $s=a+ib$ avec $a>0$.
        
            $u_n = \frac{(-1)^n}{n^s} = \frac{(-1)^n}{n^a} \cos(b\ln n) - i \frac{(-1)^n}{n^a} \sin(b\ln n)$.

            Soit $v_n = \text{Re}(u_n)$.
            $V_n = v_{2n} + v_{2n+1} = \frac{1}{(2n)^a}\cos(b\ln(2n)) - \frac{1}{(2n+1)^a}\cos(b\ln(2n+1))$.

            Développement asymptotique :
            $\frac{1}{(2n+1)^a} = \frac{1}{(2n)^a}(1+\frac{1}{2n})^{-a} = \frac{1}{(2n)^a} - \frac{a}{(2n)^{a+1}} + o(\frac{1}{n^{a+1}})$
            
            $\cos(b\ln(2n+1)) = \cos(b\ln(2n) + b\ln(1+\frac{1}{2n})) \approx \cos(b\ln(2n)) - \sin(b\ln(2n))\frac{b}{2n}$.
            
            En combinant, on obtient $V_n = O(\frac{1}{n^{1+a}})$.
            Comme $1+a > 1$, $\sum V_n$ converge absolument.

        \item 
            Puisque $v_n \to 0$, la convergence de $\sum (v_{2n}+v_{2n+1})$ implique celle de $\sum v_n$.

        \item
            Idem pour la partie imaginaire. Donc $A(s)$ converge pour $\text{Re}(s)>0$.
    \end{enumerate}
\end{enumerate}
\end{correction}

% --- Correction Exercice 4 ---
\begin{correction}{4}
\begin{enumerate}
    \item
    \begin{enumerate}
        \item $\displaystyle \sum_{k=0}^N (-x^2)^k = \frac{1-(-x^2)^{N+1}}{1-(-x^2)} = \frac{1-(-1)^{N+1}x^{2N+2}}{1+x^2}$.
        \item 
            En intégrant sur $[0,1]$ :
            $\displaystyle \int_0^1 \sum (-x^2)^k \dx = \sum_{k=0}^N (-1)^k \int_0^1 x^{2k} \dx = \sum_{k=0}^N \frac{(-1)^k}{2k+1}$.

            D'autre part : $\displaystyle\int_0^1 \frac{\dx}{1+x^2} - \int_0^1 \frac{(-1)^{N+1}x^{2N+2}}{1+x^2} \dx = [\arctan x]_0^1 - R_N = \frac{\pi}{4} - R_N$.

            Majoration du reste : $\displaystyle |R_N| \leqslant \int_0^1 x^{2N+2} \dx = \frac{1}{2N+3} \to 0$.
            Donc $\displaystyle \sum_{n=0}^{+\infty} \frac{(-1)^n}{2n+1} = \frac{\pi}{4}$.
    \end{enumerate}
    
    \item
    \begin{enumerate}
        \item 
            $\displaystyle I_{n+1}-I_n = \int_0^{\pi/2} \frac{x^2}{\sin x} (\sin((2n+2)x) - \sin(2nx)) \dx$.
            
            Utilisation de $\sin p - \sin q = 2 \sin\frac{p-q}{2} \cos\frac{p+q}{2}$.
            On a $\displaystyle I_{n+1}-I_n = \int_0^{\pi/2} 2x^2 \cos((2n+1)x) \dx$.

            Par double IPP ou calcul direct, on obtient le résultat donné.

        \item 
            $f(x) = \frac{x^2}{\sin x} \sim_0 x$. $f(0)=0$. $f$ est continue en 0.
            $f'(x) = \frac{2x\sin x - x^2\cos x}{\sin^2 x}$ et tend vers $1$ en $0$.

            Donc par prolongement $\mathscr C^1$, $f$ est $\mathscr C^1$ sur $[0,1]$.

        \item
            Par IPP sur $I_n$ (avec $u=f(x)$ et $v'=\sin(2nx)$), on obtient la majoration en $1/n$ (ou $1/2n$).
            $\sum (I_{n+1}-I_n)$ est une série télescopique qui converge vers $-I_0 = 0$ (car $I_n \to 0$ par Riemann-Lebesgue).

            La somme des termes en $1/(2n+1)$ donne $\pi^2/8$ (pré-requis ?). La somme des termes en $1/(2n+1)^3$ s'en déduit.
        
        \item La correction indique que $\sum \frac{(-1)^n}{(2n+1)^3} = \frac{\pi^3}{32}$.
    \end{enumerate}
\end{enumerate}
\end{correction}

% --- Correction Exercice 6 ---
\begin{correction}{6}
\begin{itemize}
    \item $\sum \frac{(-1)^n}{3n-1}$ : Converge par le TSSA ($1/(3n-1) \searrow 0$). Ne converge pas absolument.
    
    \item $\sum \frac{(-1)^n}{n^\alpha - (-1)^n n^\beta}$ avec $\alpha \neq \beta$.
        Supposons $\alpha > \beta$.
        $u_n = \frac{(-1)^n}{n^\alpha(1 - (-1)^n n^{\beta-\alpha})} = \frac{(-1)^n}{n^\alpha} + \frac{1}{n^{2\alpha-\beta}} + o(\dots)$. \\
        Le terme $\frac{(-1)^n}{n^\alpha}$ est le terme général d'une série convergente (TSSA). \\
        Le terme $\frac{1}{n^{2\alpha-\beta}}$ est de signe constant positif. \\
        La série converge ssi $\sum \frac{1}{n^{2\alpha-\beta}}$ converge, soit $2\alpha - \beta > 1$.

        Si $\alpha < \beta$.
        Le terme général est équivalent à $- \frac 1 {n^{\beta}}$ donc la série converge ssi $\beta > 1$ (série de Riemann).
        
    \item $\sum \ln(1 + \frac{(-1)^n (\ln n)}{n^\alpha})$.
        DL : $u_n = \frac{(-1)^n \ln n}{n^\alpha} - \frac{\ln^2 n}{2 n^{2\alpha}} + \dots$ \\
        Le premier terme converge (TSSA). Le second terme est de signe constant. \\
        La série converge ssi $2\alpha > 1$, donc $\alpha > 1/2$.
\end{itemize}
\end{correction}

\begin{correction}{7}
Soit $n \in \N$, on a la décomposition suivante :
    \[n! e = P_n + R_n  \quad \text{où} \quad P_n = \sum_{k=0}^{n} \frac{n!}{k!} \in \N \quad \text{et} \quad R_n = \frac{1}{n+1} + \sum_{k=n+2}^{+\infty} \frac{n!}{k!} = \frac{1}{n+1} + O\left(\frac{1}{n^2}\right) \] 
    
    Et donc: 
    \[ \sin(n! e 2\pi) = \sin(2\pi P_n + 2\pi R_n) \quad \text{et} \quad \sin(n! e \pi) = \sin(\pi P_n + \pi R_n) \]
    
    Puisque $P_n \in \N$, on a :
    \[ \sin(n! e 2\pi) = \sin(2\pi R_n) \quad \text{et} \quad \sin(n! e \pi) = (-1)^{P_n}\sin(\pi R_n) \]
    
    Or :
    \[\sin(2\pi R_n) \sim 2\pi R_n \sim \frac{2\pi}{n+1} \] 
    
    Par comparaison avec la série harmonique, la série de terme général $\sin(n! e 2\pi)$ diverge. \newline
    
    Mais par récurrence sur $n$ on étudie la parité de $P_n$:
    \begin{itemize}
    \item \text{Initialisation :} 
        Pour $n = 0$, $P_0 = \frac{0!}{0!} = 1$ (impair). 
        Pour $n = 1$, $P_1 = \frac{1!}{0!} + \frac{1!}{1!} = 1 + 1 = 2$ (pair).
    
    \item \text{Hérédité :} Soit $n \in \N$. Par hypothèse de récurrence, on sait que : \\
        $n \equiv 0 \ [2] \implies P_n \equiv 1 \ [2]$ \quad et \quad $n \equiv 1 \ [2] \implies P_n \equiv 0 \ [2]$ \quad Et on a \quad $P_{n+1} = (n+1)P_n + 1$ \\

        Donc, pour n pair, $P_{n+1} = \underbrace{ \underbrace{(n+1)}_{\text{impair}} \underbrace{P_n}_{\text{impair}} + 1}_{\text{pair}}$ 
        et pour n impair, $P_{n+1} = \underbrace{ \underbrace{(n+1)}_{\text{pair}} \underbrace{P_n}_{\text{pair}} + 1}_{\text{impair}}$ 
    \end{itemize}
    Ainsi, selon le principe de récurrence, on en déduit que: $\forall n \in \N, (-1)^{P_n} = (-1)^{n+1}$ \\
    
    Et par le critère des séries spéciales alternées, on en déduit que la série de terme général $\sin(n! \pi e)$ est semi-convergente.    
\end{correction}

\begin{correction}{8}
\begin{enumerate}
    \item ...
    \item On applique Cesàro.
\end{enumerate}
\end{correction}

\begin{correction}{9}
    Notons $S$ la série. On a, pour tout $n\in\N$,
    \[
    0\leq \frac{1}{(3n)!}\leq \frac{1}{n!}.
    \]
    Or, la série exponentielle converge, donc $S$ converge.
    
    Mais pour $j=e^{2i\pi/3}$, on a,
    \[
    e+e^j+e^{j^2}
    =\sum_{k=0}^{+\infty}\frac{1+j^k+j^{2k}}{k!}.
    \]
    
    Or, pour tout $k\in\N$,
    \[
    1+j^k+j^{2k}=
    \begin{cases}
        3 & \text{si } 3\mid k,\\
        0 & \text{sinon}.
    \end{cases}
    \]
    Ainsi,
    \[
    e+e^j+e^{j^2}
    =3\sum_{n=0}^{+\infty}\frac{1}{(3n)!}.
    \]
    
    Puisque $ e^j=e^{-\frac12}e^{i\frac{\sqrt{3}}{2}}
    \quad\text{et}\quad
    e^{j^2}=e^{-\frac12}e^{-i\frac{\sqrt{3}}{2}}.
    \quad\text{et par la relation}\quad
    e^{i\theta}+e^{-i\theta} = 2 \cos \theta
    $.
     
    On a :
    \[
    S = \frac{e+e^j+e^{j^2}}{3}
    = \frac13\left(e+\frac{2}{\sqrt{e}}\cos\left(\frac{\sqrt{3}}{2}\right)\right)
    \]
\end{correction}

\begin{correction}{10}
\begin{enumerate}
    \item
    \begin{enumerate}
        \item $\|u\| < 1$. La série $\sum \|u\|^n$ converge (géométrique). \\
            Donc $\sum u^n$ converge absolument (donc converge car l'espace est complet - dimension finie).
        
        \item $(e-u) \sum_{k=0}^n u^k = e - u^{n+1}$.
            Comme $u^n \to 0$, le produit tend vers $e$. \\
            Donc $(e-u)$ est inversible d'inverse la somme de la série.
    \end{enumerate}
\end{enumerate}
\end{correction}

\begin{correction}{11}
\begin{enumerate}
    \item $\|\cdot\|$ est définie positive.
    Soit $A,B \in \mathscr M_n(\C)$.

    Séparation : $\|0\| = 0$. Supposons que $\|A\| = 0$.

    Alors $\forall i,j \in \ninter 1 n,~|a_{i,j}| \leqslant 0$. Donc $A = 0$.
    
    Homogénéité : soit $\lambda \in \C$,
    \[\|\lambda A\| = \max_{i,j \in \ninter 1 n} |\lambda a_{i,j}| = \lambda \max_{i,j \in \ninter 1 n} |a_{i,j}| = \lambda \|A\|\]

    Inégalité triangulaire :
    \[\|A + B\| = \max_{i,j \in \ninter 1 n} |a_{i,j} + b_{i,j}| \leqslant \max_{i,j \in \ninter 1 n} (|a_{i,j}| + |b_{i,j}|) \leqslant \|A\| + \|B\|\]

    \item Soient $A,B \in \mathscr M_n(\C)$,
    \begin{align*}
        \|AB\| &= \max_{i,j \in \ninter 1 n} \left|\sum_{k=1}^n a_{i,k} b_{k,j}\right| \\
               &\leqslant \max_{i,j \in \ninter 1 n} \sum_{k=1}^n |a_{i,k} b_{k,j}| \\
               &\leqslant \max_{i,j \in \ninter 1 n}\left(n \max_{k \in \ninter 1 n} |a_{i,k}| |b_{k,j}|\right) \\
               &\leqslant n \left(\max_{i,k \in \ninter 1 n} |a_{i,k}|\right) \left(\max_{k,j \in \ninter 1 n} |b_{k,j}|\right) \\
               &\leqslant n \|A\| \|B\|
    \end{align*}

    Par réccurrence immédiate en prenant $B = A$, on obtient que $\forall p \in \N^*,~\|A^p\| \leqslant n^{p-1}\|A\|^p$.

    \item Soit $p \in \N^*$,
    $\frac {\|A^p\|} {p!} \leqslant \frac {n^{p-1}\|A\|^p} {p!} = \frac {(n\|A\|)^p} {p!n}$
    par critère de comparaison, le terme général est en $O\left(\frac 1 {n^2}\right)$ donc la série converge absolument.

    Étant en dimension finie, la série est convergente.
\end{enumerate}
\end{correction}

\begin{correction}{12}
\begin{itemize}
    \item $A = \begin{pmatrix} 0 & a \\ -a & 0 \end{pmatrix}$.
        $A^2 = -a^2 I_2$ et $A^3 = -a^3 J$ (avec $J=\begin{pmatrix} 0 & 1 \\ -1 & 0 \end{pmatrix}$).
        En fait $A = a J$. $J^2 = -I_2$.
            
        $\exp(A) = \sum \frac{A^k}{k!} = \sum (-1)^p \frac{a^{2p}}{(2p)!} I_2 + \sum (-1)^p \frac{a^{2p+1}}{(2p+1)!} J$.

        $\exp(A) = (\cos a) I_2 + (\sin a) J = \begin{pmatrix} \cos a & \sin a \\ -\sin a & \cos a \end{pmatrix}$.
    \item $A = \begin{pmatrix} 1 & 1 \\ 0 & 1 \end{pmatrix} = I_2 + N$ avec $N=\begin{pmatrix} 0 & 1 \\ 0 & 0 \end{pmatrix}$.
        $I_2$ et $N$ commutent. 
        
        $\exp(A) = \exp(I_2)\exp(N) = e (I_2+N) = \begin{pmatrix} e & e \\ 0 & e \end{pmatrix}$.
\end{itemize}
\end{correction}

\begin{correction}{13}
\begin{enumerate}
    \item $A = \begin{pmatrix} 3 & 6 \\ 4 & 1 \end{pmatrix}$. Polynôme caractéristique : $X^2 - 4X - 21 = (X-7)(X+3)$.
        Valeurs propres $7$ et $-3$.

        $A^n = \alpha_n A + \beta_n I$. Ou système avec les valeurs propres.
        $7^n = 7a_n + b_n$ (système adapté à la décomposition de l'énoncé).

        On trouve $A^n = \frac{1}{10} [ 7^n (A+3I) - (-3)^n (A-7I)]$.

    \item On a immédiatement $\exp(A) = \frac{e^7}{10}(A+3I) - \frac{e^{-3}}{10}(A-7I)$.
\end{enumerate}
\end{correction}

\begin{correction}{15}
\begin{enumerate}
    \item $a_{n,k} = \frac{1}{k!} \left(1 - \frac{n(n-1)\dots(n-k+1)}{n^k} \right)$.
        Or $\frac{n}{n} \dots \frac{n-k+1}{n} \leqslant 1$. Donc $a_{n,k} \ge 0$.
    
    \item
    \begin{enumerate}   
        \item $\displaystyle M_n = \left(I + \frac A n \right)^n = \sum_{k=0}^n \binom{n}{k} \frac{A^k}{n^k} = \sum_{k=0}^n \frac{1}{k!} \frac{n!}{n^k(n-k)!} A^k$.

            Différence : $\displaystyle D_n = - \sum_{k=0}^n a_{n,k} A^k$
            En norme : $\displaystyle \|D_n\| \leqslant \sum_{k=0}^n a_{n,k} \|A\|^k = \sum_{k=1}^n a_{n,k} ||A||^k$.

            Refaire le calcul avec $||A||$ revient au même donc $\displaystyle||D_n|| \leqslant \left(1 + \frac {||A||} n\right)^n - \sum_{k=0}^n \frac {||A||^k} {k!}$.

            On utilise l'exponentielle réelle : $\left(1+\frac{\|A\|}{n}\right)^n \to e^{\|A\|}$.
            L'autre aussi, donc la différence tend vers 0.

        \item On en déduit directement que $M_n \to \exp(A)$.
    \end{enumerate}
\end{enumerate}
\end{correction}

\begin{correction}{17}

    $|u_n| \sim \frac 1 {n^2}$. Donc la série converge.

    Calculons les sommes partielles :
    $$\sum_{k=2}^n \ln(k^2 - 1) - \ln(k^2) = \sum_{k=3}^{n+1} \ln(k) + \sum_{k=1}^{n-1} \ln(k) - 2 \sum_{k=2}^{n} \ln(k) = \ln(1 + \frac 1 n) - \ln(2)$$
    Donc la somme est $-\ln(2)$.
\end{correction}

\begin{correction}{19}
\begin{enumerate}
    \item   
        $u_n > 0$. $\frac{u_{n+1}}{u_n} = \frac{(n+a)(n+b)}{(n+c)(n+d)} \to 1$. Cas douteux d'Alembert.

        $v_n = \ln(n^\alpha u_n)$.
        $v_{n+1}-v_n = \alpha \ln(1+\frac{1}{n}) + \ln \frac{u_{n+1}}{u_n}$.

        DL de $\frac{u_{n+1}}{u_n} = 1 + \frac{a+b-c-d}{n} + O\left(\frac{1}{n^2}\right)$.
        $v_{n+1}-v_n = \frac{\alpha + a+b-c-d}{n} + O\left(\frac{1}{n^2}\right)$.

        Série télescopique converge ssi le terme en $1/n$ est nul.
        Condition : $\alpha = c+d-a-b$.

    \item Si cette condition est remplie, $v_n \to L$.
        Donc $n^\alpha u_n \to e^L$,
        donc $u_n \sim \frac{K}{n^{c+d-a-b}}$.
\end{enumerate}
\end{correction}

\begin{correction}{22}
\begin{enumerate}
    \item Si $\frac{u_{n+1}}{u_n} \leqslant \frac{v_{n+1}}{v_n}$, alors $\frac{u_{n+1}}{v_{n+1}} \leqslant \frac{u_n}{v_n}$.
        La suite $(u_n/v_n)$ est décroissante, donc bornée.
        
        $u_n = O(v_n)$. Si $\sum v_n$ converge, alors $\sum u_n$ converge.
        
    \item Comparaison avec Riemann $v_n = \frac{1}{n}^\beta$.

        $\frac{v_{n+1}}{v_n} = \left(1+\frac{1}{n}\right)^{-\beta} = 1 - \frac{\beta}{n} + o\left(\frac{1}{n}\right)$.
        Si $\frac{u_{n+1}}{u_n} = 1 - \frac{\alpha}{n} + o(\frac{1}{n})$.
    \begin{itemize}
        \item Si $\alpha > 1$, on peut choisir $\beta$ tel que $1 < \beta < \alpha$.

            Alors pour $n$ grand, $1 - \frac{\alpha}{n} < 1 - \frac{\beta}{n}$, donc $\frac{u_{n+1}}{u_n} < \frac{v_{n+1}}{v_n}$.
            $\sum v_n$ converge (Riemann $\beta>1$), donc $\sum u_n$ converge.

        \item Si $\alpha < 1$, divergence par comparaison avec $\beta \in ]\alpha, 1[$.
    \end{itemize}
\end{enumerate}
\end{correction}

\begin{correction}{23}
\begin{itemize}
    \item $u_n = \frac 1 {n^{\ln(\ln(n))}} \leqslant \frac 1 {n^2}$ pour $n$ assez grand donc elle converge. 
    \item $u_n = \frac{n^n}{e^{n^2}}$. $n^2 u_n \to 0$. Donc elle converge (racine n-ième).
    \item $u_n = \frac{n!^2}{(2n)!}$. Stirling : $\frac{n^{2n} e^{-2n} (2\pi n)}{(2n)^{2n} e^{-2n} \sqrt{4\pi n}} \sim \frac{\sqrt{\pi n}}{4^n}$. D'Alembert. Converge

    \item $u_n = (\sin \frac{1}{n}) n^\alpha \sim n^{\alpha-1}$. Converge ssi $1-\alpha > 1 \iff \alpha < 0$. 
    
        (Attention, énoncé dit $e^{n^\alpha \ln(...)}$, vérifier l'énoncé exact, correction mentionne exponentielle).
        L'énoncé PDF source mentionne $\sum (\sin \frac{1}{n}) n^\alpha$.
        La correction manuscrite traite $u_n = e^{n^\alpha \ln(n \sin(1/n))}$.

        $n \sin(1/n) = 1 - \frac{1}{6n^2}$.
        $\ln(\dots) \sim -\frac{1}{6n^2}$.

        Exposant $\sim - \frac{n^{\alpha-2}}{6}$.
        \begin{itemize}
            \item Si $\alpha < 2$, exposant $\to 0$, $u_n \to 1$. Diverge.
            \item Si $\alpha = 2$, $u_n \to e^{-1/6}$. Diverge.
            \item Si $\alpha > 2$, exposant $\to -\infty$ (très vite ? ou vers 0 par valeurs négatives ?).
        \end{itemize}
        Attention, il y a $n^\alpha$ en facteur devant le log dans l'exposant.

        Si $\alpha > 2$, $n^{\alpha-2} \to +\infty$, terme $\to e^{-\infty} = 0$.
        Règle $n^2 u_n$.
\end{itemize}
\end{correction}

\begin{correction}{24}
\begin{enumerate}
    \item Si $\sum u_n$ converge, $S_n \to S > 0$. Donc $u_n/S_n \sim u_n/S$. Converge.
    \item Si $\sum u_n/S_n$ converge.
        $\ln(\frac{S_n}{S_{n-1}}) = \ln(1 + \frac{u_n}{S_{n-1}}) \approx \frac{u_n}{S_{n-1}} \approx \frac{u_n}{S_n}$.

        La série $\sum \ln(S_n) - \ln(S_{n-1})$ est de même nature que $\sum u_n/S_n$.
        C'est une somme télescopique $\ln S_n$.

        Si elle converge, $\ln S_n$ a une limite, donc $S_n$ a une limite finie. $\sum u_n$ converge.
\end{enumerate}
\end{correction}

\begin{correction}{25}
\begin{enumerate}
    \item 
        $\displaystyle R_n = \sum_{k=n+1}^{+\infty} \frac{1}{k^\alpha}$.
        Comparaison série-intégrale : $\displaystyle R_n \sim \int_n^{+\infty} \frac{\dt}{t^\alpha} = \frac{1}{(\alpha-1)n^{\alpha-1}}$.
        
    \item $\sum R_n$ converge ssi $\sum \frac{1}{n^{\alpha-1}}$ converge
        ssi $\alpha-1 > 1 \iff \alpha > 2$.

    \item Somme : 
        \begin{align*}
            \sum_{n=1}^{+\infty} R_n &= \sum_{n=1}^{+\infty} \sum_{k=n+1}^{+\infty} \frac{1}{k^\alpha} \\
                &= \sum_{k=2}^{+\infty} \frac{1}{k^\alpha} \sum_{n=1}^{k-1} 1 \\
                &= \sum_{k=2}^{+\infty} \frac{k-1}{k^\alpha} \\
                &= \zeta(\alpha-1) - \zeta(\alpha) &&
        \end{align*}
\end{enumerate}
\end{correction}


